\documentclass[12pt]{article}
\usepackage[T1]{fontenc}
\usepackage[utf8]{inputenc}
\usepackage{lmodern}
\usepackage{textcomp}
\usepackage{enumitem}
\usepackage{geometry}
\geometry{margin=1in}
\begin{document}
\begin{center}
Answer Key - Economics - K-12 Level Assessment 25 \
\small (Answers are ordered exactly as in the question paper.)
\end{center}

\section*{Multiple Choice}
\begin{enumerate}[label=\arabic*.]
\item \textbf{Answer:} D
\\ \textit{Options:} A) less choices due to inflation.; B) financial assets being worth less due to inflation.; C) a la carte savings falling.; D) resource misallocation due to inflation.
\\ \textit{Note:} The correct answer is accurate because "menu costs" describe the expenses incurred by businesses in changing prices, which can lead to inefficient allocation of resources in the economy during periods of inflation.
\item \textbf{Answer:} B
\\ \textit{Options:} A) The FED lowers the discount rate.; B) The FED sells government securities in the secondary market.; C) The federal government spends less money.; D) The FED lowers reserve requirements.
\\ \textit{Note:} The correct answer is that when the Federal Reserve sells government securities in the secondary market, it reduces the amount of money in circulation because buyers pay for these securities, which takes money out of the economy.
\item \textbf{Answer:} D
\\ \textit{Options:} A) decreases over time; B) equals the full-employment level of real GDP; C) is above the full-employment level of real GDP; D) is below the full-employment level of real GDP
\\ \textit{Note:} A recessionary gap occurs when the actual output of the economy, represented by the short-run equilibrium level of real GDP, falls short of the potential output, or full-employment level of real GDP, indicating that resources are underutilized.
\item \textbf{Answer:} C
\\ \textit{Options:} A) The equilibrium price will rise and the equilibrium quantity will fall.; B) The equilibrium price will fall and the equilibrium quantity will rise.; C) The equilibrium price and the equilibrium quantity will both rise.; D) The equilibrium price and the equilibrium quantity will both fall.
\\ \textit{Note:} When the price of good X, a close substitute for good Z, rises, consumers will shift their demand from good X to good Z, leading to an increase in both the equilibrium price and quantity of good Z.
\item \textbf{Answer:} C
\\ \textit{Options:} A) As the demand for bonds increases the interest rate increases.; B) For a given money supply if nominal GDP increases the velocity of money decreases.; C) When demand for stocks and bonds increases the asset demand for money falls.; D) A macroeconomic recession increases the demand for loanable funds.
\\ \textit{Note:} As demand for stocks and bonds increases, investors prefer to allocate their funds into these assets rather than holding onto cash, resulting in a decrease in the asset demand for money.
\end{enumerate}

\section*{True / False}
\begin{enumerate}[label=\arabic*.]
\setcounter{enumi}{5}
\item \textbf{Answer:} False
\\ \textit{Options:} A) True; B) False
\\ \textit{Note:} Flash estimates of GDP are often preliminary figures that are subject to revisions as more comprehensive data becomes available, thus they do require updates after publication.
\item \textbf{Answer:} True
\\ \textit{Options:} A) True; B) False
\\ \textit{Note:} The statement is true because M$_{1}$ includes the most liquid forms of money, such as cash and demand deposits, while M$_{2}$ encompasses additional assets like savings accounts that, although can be converted into cash, are not as readily accessible for immediate transactions.
\item \textbf{Answer:} True
\\ \textit{Options:} A) True; B) False
\\ \textit{Note:} Keynesian economists argue that during periods of economic downturns, government spending and fiscal policy can more directly increase demand and stimulate economic activity than monetary policy, which may be less effective in such situations.
\item \textbf{Answer:} True
\\ \textit{Options:} A) True; B) False
\\ \textit{Note:} When the price of one good increases, consumers tend to buy more of a substitute good instead, leading to an increase in its quantity demanded, which defines them as substitute goods.
\item \textbf{Answer:} False
\\ \textit{Options:} A) True; B) False
\\ \textit{Note:} A tax cut in a full employment economy can increase disposable income, leading to higher consumer spending and potentially raising real output rather than lowering it, while the effects on the price level depend on various factors like aggregate demand and supply.
\end{enumerate}

\section*{Long Form Response}
\begin{enumerate}[label=\arabic*.]
\setcounter{enumi}{10}
\item \textbf{Answer:} Labor unions utilize various methods to increase wages for their members, including collective bargaining, strikes, and lobbying for favorable legislation. One effective strategy is collective bargaining, where unions negotiate directly with employers to secure better pay and benefits. Additionally, unions may work to decrease the prices of substitute resources, such as promoting the use of skilled labor over cheaper alternatives, which can create a stronger demand for their members' work. However, some approaches, like aggressive strikes, can be risky and may lead to job losses or negative public perception, making them less frequently used. Overall, strategies that focus on collaboration and skill enhancement tend to be more effective and sustainable in achieving wage increases.
\\ \textit{Note:} correct because it identifies key strategies employed by labor unions, such as collective bargaining and lobbying, to effectively negotiate higher wages while also acknowledging the importance of promoting skilled labor to enhance demand for union members' work.
\item \textbf{Answer:} A rightward shift in the aggregate demand curve leads to an increase in both the equilibrium price level and the quantity of output in the economy. This shift indicates that consumers, businesses, and the government are collectively demanding more goods and services at every price level. Since the aggregate supply curve is upward sloping, producers respond to the increased demand by raising prices and increasing production to meet the higher demand. As a result, the overall price level rises, and the total output in the economy also increases, reflecting a more active economic environment.
\\ \textit{Note:} correct because a rightward shift in the aggregate demand curve signifies increased demand across the economy, which, in conjunction with an upward-sloping aggregate supply curve, results in higher prices and greater output as producers respond to the heightened demand.
\end{enumerate}

\vfill
\noindent\textit{End of Answer Key}
\end{document}